\documentclass[11pt,a4paper]{article}

% ---------- Encodage et langue ----------
\usepackage[utf8]{inputenc}
\usepackage[T1]{fontenc}
\usepackage[french]{babel}
\usepackage{lmodern,microtype}

% ---------- Mise en page ----------
\usepackage[
    left=2.2cm,
    right=2.2cm,
    top=1.7cm,
    bottom=1.7cm
]{geometry}

\usepackage[hidelinks]{hyperref}

% Faciliter l'extraction du texte du PDF
\input{glyphtounicode}
\pdfgentounicode=1

% ---------- Champs à compléter ----------
\newcommand{\Entreprise}{[Nom de l'entreprise]}
\newcommand{\AdresseEntreprise}{[Adresse de l'entreprise]}
\newcommand{\DateEnvoi}{[Date d'envoi]}

\pagestyle{empty}

\setlength{\parindent}{0pt}
\setlength{\parskip}{0.65em}

\begin{document}

% ============================================================
% EXPÉDITEUR
% ============================================================

\textbf{Céleste KOLOUSSA}\\
Nantes / Saint-Nazaire\\
07 83 60 11 87\\
\href{mailto:ckoloussa@gmail.com}{ckoloussa@gmail.com}\\
\href{https://celestekoloussa.fr}{celestekoloussa.fr}

% ============================================================
% DESTINATAIRE ET DATE
% ============================================================

%\begin{flushright}
%    \Entreprise\\
%    À l'attention du service recrutement\\
%    \AdresseEntreprise

%    Nantes, le \DateEnvoi
%\end{flushright}

\vspace{1cm}
% ============================================================
% OBJET
% ============================================================

\textbf{Objet : Candidature spontanée ---
technicien en électronique et systèmes embarqués}

\vspace{0.5cm}
% ============================================================
% LETTRE
% ============================================================

Madame, Monsieur,

Titulaire d'un BUT Génie électrique et informatique
industrielle, parcours Électronique et systèmes embarqués,
je vous propose ma candidature pour un poste de technicien.
Je souhaite travailler sur la conception, les essais ou
le diagnostic de systèmes électroniques, en associant
les mesures sur le matériel au développement logiciel.
Mon projet de poursuite d'études n'ayant pas pu se concrétiser
faute d'alternance, je souhaite aujourd'hui m'engager dans
un emploi technique et progresser au contact d'une équipe.

Chez LABAA, qui développe des fauteuils thérapeutiques,
j'ai travaillé sur des problèmes de régulation de cartes
électroniques. J'ai effectué des mesures en conditions
d'utilisation, recherché des composants de remplacement
et participé à leur validation. J'ai aussi pris en charge
des cartes retournées au SAV, du diagnostic au rapport
de panne et au remplacement des composants. Pour faciliter
leur programmation, j'ai développé un outil sur Raspberry Pi,
en Python avec PlatformIO, commandé par un bouton et muni
d'une LED indiquant son état.

Au laboratoire SUBATECH, dans le cadre du projet de caméra
médicale XEMIS2, j'ai posé les premières bases d'un logiciel
d'acquisition en C++ et Qt sous Linux. Il permettait
de récupérer des données depuis un serveur, de les afficher
en temps réel, puis de les enregistrer dans un format binaire
pour les relire avec leurs repères temporels. Ce travail
m'a amené à relier les besoins de visualisation des scientifiques
aux choix de structure du logiciel et de stockage des données.

Mon engagement comme vice-président de l'association
robotique de l'Université de Nantes m'a appris à construire
un projet avec une équipe : répartir les fonctions, choisir
les composants et faire avancer ensemble l'électronique,
le firmware et la mécanique. Dans un projet universitaire
sous ROS1, j'ai notamment ajouté une carte d'interfaçage ESP32
et des capteurs ultrasonores pour compléter la détection
du LiDAR et arrêter le robot à l'approche d'un obstacle,
en conduite manuelle comme autonome.

Je pourrais ainsi contribuer à vos activités de bureau
d'études, d'essais, de maintenance ou de production
électronique. J'aime comprendre l'origine d'un défaut,
tester une solution et expliquer ce qui a été réalisé.
Titulaire du permis B et véhiculé, je suis disponible
pour échanger sur vos besoins et les missions que
vous pourriez me confier.

Je vous prie d'agréer, Madame, Monsieur, l'expression
de mes salutations distinguées.

\hfill\textbf{Céleste KOLOUSSA}

\end{document}